% 图 1.1 系统概述图（重绘加强版）
% 基于 CICC0903540 技术文档 1.1 - 语义通信图像传输系统
\documentclass[tikz,border=10pt]{standalone}

% 强制全局样式与配色
\usepackage{amssymb, amsmath}
\usepackage{fontspec}
\usepackage{xeCJK}
\setCJKmainfont{SimHei}
\usepackage{tikz}
\usetikzlibrary{
  arrows.meta, positioning, shapes.geometric, calc,
  backgrounds, fit, decorations.pathreplacing, patterns
}

% 低饱和度马卡龙配色
\definecolor{mBlue}{HTML}{AEC6CF}
\definecolor{mPink}{HTML}{FFD1DC}
\definecolor{mGreen}{HTML}{B1D8B7}
\definecolor{mPurple}{HTML}{B39EB5}
\definecolor{mYellow}{HTML}{FDFD96}
\definecolor{mGray}{HTML}{E5E4E2}
\definecolor{mDark}{HTML}{4A4A4A}

% 全局 TikZ 样式
\tikzset{
  >=stealth,
  thick,
  draw=mDark,
  block/.style={
    rectangle, rounded corners=3pt, draw=mDark, align=center,
    font=\sffamily\small, minimum height=0.95cm, inner sep=6pt},
  micro/.style={
    rectangle, rounded corners=2pt, draw=mDark!70, align=center,
    font=\sffamily\tiny, minimum height=0.35cm, inner sep=2pt},
  ts/.style={font=\tiny\sffamily, color=mDark!80},
  arr/.style={-{Stealth[length=4pt]}, thick, draw=mDark},
  legendbox/.style={
    rectangle, draw=mDark!50, rounded corners=3pt,
    fill=mGray!15, inner sep=5pt, font=\tiny\sffamily},
}

\begin{document}
\begin{tikzpicture}[
  node distance=0.6cm and 1.5cm,
]
  % ========== 发送端（上位机）==========
  \node[font=\sffamily\footnotesize\bfseries, color=mDark] (txLabel) {发送端};

  \node[block, fill=mBlue!40, below=0.4cm of txLabel, minimum width=2.8cm] (enc)
    {Encoder\\{\footnotesize 语义特征提取}};

  \node[block, fill=mBlue!40, below=0.5cm of enc, minimum width=2.8cm] (quant)
    {量化压缩\\{\footnotesize Quantization}};

  % Encoder 微观结构
  \node[micro, fill=mBlue!30, right=0.2cm of enc, xshift=0.5cm] (e1) {Conv};
  \node[micro, fill=mBlue!30, right=0.12cm of e1] (e2) {BN};
  \node[micro, fill=mBlue!30, right=0.12cm of e2] (e3) {ReLU};
  \node[micro, fill=mBlue!25, right=0.15cm of e3] (e4) {Res$\times$2};
  \draw[arr, -, thin] (e1) -- (e2) -- (e3) -- (e4);
  \draw[arr, dashed, thin, mDark!60] (enc.east) |- (e2.west);

  % 硬件配置
  \node[block, fill=mGray!60, below=0.4cm of quant, minimum width=2.8cm,
    font=\sffamily\scriptsize] (pc)
    {上位机（客户端）\\i7-13700H + RTX 4060\\8GB GPU + 16GB RAM};

  % ========== 传输信道 ==========
  \node[font=\sffamily\footnotesize\bfseries, color=mDark, right=5.0cm of txLabel] (chLabel)
    {传输信道};

  \node[block, fill=mGreen!50, below=0.4cm of chLabel, minimum width=2.5cm] (ssh)
    {OpenSSH\\{\footnotesize 安全传输通道}};

  % SSH 细节
  \node[micro, fill=mGreen!40, below=0.3cm of ssh, xshift=-0.6cm] (aes) {AES-128-CTR};
  \node[micro, fill=mGreen!40, right=0.1cm of aes] (sftp) {SFTP Pool};
  \node[micro, fill=mGreen!40, below=0.25cm of aes] (win) {Window 1GB};
  \node[micro, fill=mGreen!40, right=0.1cm of win] (chunk) {Chunk 64MB};

  % ========== 接收端（飞腾派）==========
  \node[font=\sffamily\footnotesize\bfseries, color=mDark, right=5.0cm of chLabel] (rxLabel)
    {接收端};

  \node[block, fill=mPurple!40, below=0.4cm of rxLabel, minimum width=2.8cm] (gen)
    {Generator\\{\footnotesize 图像解码与重建}};

  \node[block, fill=mPurple!40, below=0.5cm of gen, minimum width=2.8cm] (opt)
    {推理加速\\{\footnotesize TVM / MNN}};

  % Generator 微观结构
  \node[micro, fill=mPurple!30, right=0.2cm of gen, xshift=0.5cm] (g1) {Conv};
  \node[micro, fill=mPurple!30, right=0.12cm of g1] (g2) {BN};
  \node[micro, fill=mPurple!30, right=0.12cm of g2] (g3) {ReLU};
  \node[micro, fill=mPurple!25, right=0.15cm of g3] (g4) {Res$\times$5};
  \draw[arr, -, thin] (g1) -- (g2) -- (g3) -- (g4);
  \draw[arr, dashed, thin, mDark!60] (gen.east) |- (g2.west);

  % 硬件配置
  \node[block, fill=mGray!60, below=0.4cm of opt, minimum width=2.8cm,
    font=\sffamily\scriptsize] (ft)
    {飞腾派（服务端）\\ARMv8.0 Cortex-A72\\2$\times$1.8GHz + 2$\times$1.5GHz};

  % ========== 主数据流连线 + 张量维度标注 ==========
  \draw[arr] (enc) -- (quant)
    node[pos=0.5, left=2pt, ts] {$1{\times}32{\times}32{\times}32$}
    node[pos=0.5, right=2pt, ts, align=center] {语义\\特征};

  \draw[arr] (quant.east) -- ++(0.8,0) |- (ssh.west)
    node[pos=0.2, above=2pt, ts] {隐式张量（量化后）}
    node[pos=0.2, below=2pt, ts] {$\approx$32KB};

  \draw[arr] (ssh.east) -- ++(0.8,0) |- (gen.west)
    node[pos=0.2, above=2pt, ts] {$\hat{\mathbf{y}}$}
    node[pos=0.2, below=2pt, ts] {AWGN 扰动};

  \draw[arr] (gen) -- (opt)
    node[pos=0.5, left=2pt, ts] {$1{\times}3{\times}256{\times}256$}
    node[pos=0.5, right=2pt, ts, align=center] {重建\\图像};

  % ========== 层次化表达：背景分组 ==========
  \begin{scope}[on background layer]
    % 发送端区域
    \node[fit=(txLabel)(enc)(quant)(pc)(e1)(e4),
      fill=mBlue!15, draw=mDark!40, dashed,
      rounded corners=5pt, inner sep=10pt] {};

    % 传输信道区域
    \node[fit=(chLabel)(ssh)(aes)(chunk),
      fill=mGreen!15, draw=mDark!40, dashed,
      rounded corners=5pt, inner sep=10pt] {};

    % 接收端区域
    \node[fit=(rxLabel)(gen)(opt)(ft)(g1)(g4),
      fill=mPurple!15, draw=mDark!40, dashed,
      rounded corners=5pt, inner sep=10pt] {};
  \end{scope}

  % ========== 信息密度拉满：系统流程标注 ==========
  \node[legendbox, anchor=north west] at ([xshift=5pt, yshift=-5pt]current bounding box.north west)
    [align=left, text width=4.8cm] {
    \textbf{系统工作流程}\\[2pt]
    1. 上位机执行 \texttt{test\_sub.py} 编码\\
    2. OpenSSH 建立安全连接（Port 22）\\
    3. SFTP 传输隐式张量到飞腾派\\
    4. 远程触发 \texttt{decoder.py} 解码\\
    5. 输出 PNG 图像到指定目录
  };

  % ========== 图例（右下角）==========
  \node[legendbox, anchor=south east] at ([xshift=-5pt, yshift=5pt]current bounding box.south east)
    [align=left, text width=4.2cm] {
    \textbf{图例}\\[1pt]
    \colorbox{mBlue!40}{\rule{0.16cm}{0.12cm}\hspace{0.05cm}} 编码端 \quad
    \colorbox{mPurple!40}{\rule{0.16cm}{0.12cm}\hspace{0.05cm}} 解码端 \quad
    \colorbox{mGreen!50}{\rule{0.16cm}{0.12cm}\hspace{0.05cm}} 通道\\[1pt]
    数据集：Open Images V7（110k张）\\
    测试集：Places365（300张 256$\times$256）
  };
\end{tikzpicture}
\end{document}
